\documentclass{article}
\title{Review Response for ``Adaptive Projected Two-Sample Comparisons for Single-Cell Gene Expression Data''}
\date{\vspace{-5ex}}
\date{}  % Toggle commenting to test
\usepackage{multirow} 

\usepackage{microtype}
\usepackage{graphicx}
\usepackage{subfigure}
\usepackage{booktabs} % for professional tables
\usepackage{amsmath}
\usepackage{amsthm,amssymb}
\usepackage{amsfonts}
\usepackage{framed}

\usepackage{hyperref}
\hypersetup{
  pdfinfo={
  pdfproducer={},
  Title={},
  Subject={},
  Author={},
  }
}
\usepackage{abstract}
\usepackage{verbatim}
\usepackage{algorithm}
%\usepackage{algorithmic}
\usepackage{algpseudocode}
\usepackage{amsmath,amssymb,bbm,amsthm,mathrsfs,bm}
\usepackage{geometry}
\usepackage{xcolor}
\usepackage{booktabs}
\usepackage{graphicx}
\usepackage{caption}
\newtheorem{thm}{Theorem}
\theoremstyle{definition}
\geometry{left=2cm,right=2cm,top=2cm,bottom=2cm}
\renewcommand{\normalsize}{\fontsize{13pt}{\baselineskip}\selectfont}
\linespread{1.4}
\setlength{\parskip}{1em}
\DeclareMathOperator*{\argmin}{argmin}
\newcommand{\var}{\texttt}
\newcommand{\assign}{\leftarrow}
\newcommand{\multilinestate}[1]{%
  \parbox[t]{\linewidth}{\raggedright\hangindent=\algorithmicindent\hangafter=1
    \strut#1\strut}}
\newcommand{\vertiii}[1]{{\left\vert\kern-0.25ex\left\vert\kern-0.25ex\left\vert #1 
    \right\vert\kern-0.25ex\right\vert\kern-0.25ex\right\vert}}

\newcommand{\tz}{\color{teal} TZ: }
\newcommand{\kr}{\color{magenta} KR:\ }


\begin{document}
\definecolor{shadecolor}{gray}{0.9}
\maketitle
\vspace{-30mm}



% \textcolor{red}{
% 1. Explain the difference between cross-fitting debiasing and universal splitting inference.
% Dimension-agnostic inference using cross U-statistics, Kim \& Ramdas, 2024\\
% 2. Mention that the cross-fitting framework needs to be checked in different scenarios, and we are the first to do this for PCA/Sparse PCA\\
% 3. Remove the first theorem.  Start directly from estimating the sparse-projected difference.  If we only want a sparse projection-based test, without effect size estimation, then introduce anchored test.
% }

\newpage
\section{Reviewer 1's Comments}

\begin{shaded}
This manuscript proposes debiased and anchored projected two-sample tests for high-dimensional mean comparison, motivated by single-cell RNA-seq applications. The authors leverage sparse PC directions and one-step semiparametric corrections to obtain asymptotically normal test statistics, and introduce an “anchored projection” that combines discriminative directions (e.g., logistic Lasso) with sparse PCs. Simulation studies and a lupus scRNA-seq dataset are used to illustrate the methods. Overall, the problem is important and the paper is technically strong and well written, but I have several methodological and presentation concerns that I think should be addressed before publication.
\end{shaded}

We appreciate your time and effort in reviewing this manuscript.

\begin{shaded}
Major Comments\\
1. Clarify and sharpen the main methodological contribution.
It would help if the introduction and early parts of Section 2 more clearly delineated which components are novel versus adaptations of existing semiparametric / high-dimensional inference tools. For instance, to what extent is the one-step correction with sparse PCs conceptually new relative to previous debiasing approaches for high-dimensional linear functionals, and what is the precise innovation in the “anchored projection” compared with prior two-sample tests based on learned discriminant directions?
\end{shaded}

Thank you for the suggestions. We added more details in the introduction to highlight the contribution of this work:

\begin{itemize}
    \item For the one-step correction proposal, we work under the general theoretical framework. However, an eigenvector is not a simple linear mapping of the distribution. The exact details of the one-step correction, as well as its theoretical guarantee, were novel in this work. 
    \item For the anchored projection proposal, we emphasized that regular Gaussian statistics can be achieved under mild conditions without debiasing. Cross-fitting framework, using PC1 as the anchoring proxy and the equivalent means under the null hypothesis enabled this interesting regime. It has never been explored in the literature that this simple and application-relevant scheme would lead to regular inferential procedures. The closest one in the literature would be \cite{liu2022multiple}: However, their result was not established under the null hypothesis (because there was no proxy anchoring direction), the classifier aims at the unique optimal direction, and they implemented sample-splitting rather than cross-fitting (although they proposed procedures to enhance the power due to a smaller sample size).
\end{itemize}

\begin{shaded}
2. Position the work more explicitly in the high-dimensional two-sample testing literature.
The paper cites several related methods (e.g., tests based on global T2–type statistics, sparse discriminant directions, and clustering-based approaches), but the connections remain somewhat qualitative. A brief subsection or paragraph comparing the proposed tests to key existing two-sample tests—both in terms of assumptions (e.g., sparsity structure, covariance conditions) and the null/alternative they target—would help readers better understand when the proposed methods are preferable.
\end{shaded}

We appreciate your suggestion. We added one paragraph at the end of the introduction section to discuss several canonical high-dimensional mean comparison methods. Although this work targets high-dimensional mean comparison, the specific mathematical formulation of our problem is novel and more structured compared with the most commonly considered global null. Our method was motivated more by the debiasing and cross-validation literature---so in this paragraph, we also emphasize this difference and try not to mislead the readers to the line of research reviewed in \cite{huang2022overview}.

\begin{shaded}
    3. Interpretation and practical plausibility of Assumption 2.2 and Theorem 2.3 conditions. The conditions around approximate consistency of $u^{(-m)}$, the uniform integrability assumptions, and the $n^{-1 / 4}$ rate requirements for estimating $\mu_X$, $\mu_Z$ and $\Sigma$ in Theorem 2.3 are mathematically reasonable but may be hard to verify in practice. It would be useful to (i) provide more intuition about when these conditions are expected to hold in realistic scRNA-seq settings, and (ii) comment on whether milder rates (e.g., allowing somewhat slower convergence of $\Sigma$ ) would still yield a valid approximation in finite samples.
\end{shaded}

Thank you for the comment.

We use the uniform integrability assumption to rigorously state and establish the law of large numbers in our high-dimensional setting. They are basically stating that the sample means of some statistics would converge to their population means. In a high-dimensional setting, the samples are from distributions of ``increasing" dimension, the population means are no longer fixed quantities (see the triangular array comment in Remark~2.2), and the standard law of large numbers does not strictly apply. This is almost the minimal assumption for statistics to function, and we believe it is mild.

The $n^{-1/4}$ convergence rate for the covariance matrix is standard in the sense that difference debiasing problems usually end up requiring certain nuisance parameters to be estimated faster than this rate. However, as you noted, it is not trivial. In the spiked model setting, the minimax convergence rate is roughly $\sqrt{r\log p/n}$ where $p$ is the feature dimension and $r$ is the number of spikes \cite{cai2015optimal}. For polynomially increasing $p$ and fixed $r$, this rate is well-below the required $n^{-1/4}$.

It was not easy to fundamentally improve this condition for the debiased proposal, which is similar for many typical semiparametric methods in the literature. However, the approximately orthogonality proposal in this manuscript does function under a less demanding setting: it does not need explicit estimation of the covariance matrix, and does not rely on matrix estimation quality assumptions.

% (although it is a theoretically sufficient condition, not a necessary one, which didn't completely rule out the possibility). We found the other proposal established under approximately orthogonality is easy to engage with (although it trades for a weaker conclusion): it does not need explicit estimation of the covariance matrix, and the empirical estimation of classifier coefficients is less demanding under a finite sample setting.

\begin{shaded}
4. Choice and implementation details of sparse PCA.
Since the performance of both debiased and anchored tests depends strongly on the estimated sparse PCs, please provide more concrete information: which exact sPCA algorithm was used (e.g., reference and software), how sparsity/tuning parameters were chosen, and how sensitive the results are to these choices. A short sensitivity analysis (e.g., different numbers of nonzero loadings or alternative sPCA methods) would be helpful.
\end{shaded}

We implemented the $\ell_1$-penalized sparse PCA framework introduced by \cite{witten2009penalized}. We added the following details in the appendix of the paper. 

\begin{quote}
The authors of \cite{witten2009penalized} applied an iterative algorithm to solve a penalized matrix decomposition of the data matrix, which can be viewed as a low-rank singular value decomposition with an $\ell_1$ constraint on the singular vectors. The explicit objective function is given in Equation~(2.7) of \cite{witten2009penalized}. The tuning parameter in this optimization problem is selected by cross-validation. The details of this procedure are described in the manual for the \texttt{SPC.cv} function in the \texttt{R} package \texttt{PMA} \cite{PMA}. We include a short summary here for readers' convenience.

The \texttt{SPC.cv} function performs cross-validation for the first sparse principal component. It repeats the following steps across folds: (1) replace a fraction of the entries in the data matrix with missing values; (2) apply \texttt{SPC} to this partially observed data matrix over a range of tuning parameter values, each time obtaining a rank-one approximation $udv^\top$ with sparse $v$; and (3) measure the mean squared prediction error for the missing entries created in step 1. The selected tuning parameter is the one with the lowest average mean squared prediction error in step 3.

To perform cross-validation for the second sparse principal component, the same procedure is applied to the deflated matrix $X-udv^\top$, where $udv^\top$ is the rank-one component obtained by applying \texttt{SPC} to the original data matrix $X$.
\end{quote}

We also conducted a set of simulations with a different strategy for PC estimation, detailed in the response of question 5 below. 

\begin{shaded}  
5. Simulation design and coverage.
The simulations nicely demonstrate type-I error and power for a specific zero-inflated Gaussian block-correlation structure. Given the broad applicability of the method, I wonder if the authors could (or at least briefly discuss): i) additional covariance structures (e.g., more diffuse correlations or non-block patterns); ii)heavier-tailed distributions or different sparsity levels in the mean shift; iii) more extreme sample-size imbalance between groups. iv) Even a short summary in the Supplement (with pointers in the main text) would reassure readers about robustness.
\end{shaded}

To verify the robustness of the proposal, a set of supplementary simulations was added to the appendix section ``Simulation Under Unstructured Sparse Correlation". This study aims to provide more variety in the simulation setting and method details. Specifically: 1) We considered unstructured covariance matrices. They do not have block correlation; rather, they are row-sparse that permits effective estimation. 2) As a result, the true leading PC vectors are not sparsely supported. But it is still possible to estimate them because the covariance matrix has a tamed structure and enough eigen-gaps. 3) We applied hard-thresholding to estimate the sparse covariance matrix first and then compute the eigen systems. This strategy is supported by rigorous literature work (see detailed discussion in the appendix) and cooperates well with the proposed framework.

We also expanded the discussion in section ``Debiased Tests for the Projected Null" to address the other comments. We added a short note on near-sparse means estimation---they may have full supports but not too many entries can have large values. Estimating a high-dimensional mean without constraints in general is very challenging when the errors are measured in $\|\cdot\|_2$-norm. Sparse covariance matrix estimation is also possible using robust thresholding estimators: these methods typically only requires the distributions have entry-wise forth-moments, instead of the light-tail assumptions that require all moments.

\begin{shaded}
6. Interpretation of the lupus application results.
The real-data analysis is interesting, but the statistical conclusions could be translated more clearly into biological insight. For example, beyond reporting p-values and overlap counts, it would be helpful to more explicitly connect what the significant PCs and anchored directions suggest about specific pathways or cell-type–specific processes, and to contrast these findings with those from standard marginal methods (e.g., what new insight do we gain by using projected tests rather than a list of gene-level p-values?).
\end{shaded}

Thank you for the question, and it prompted us to think more carefully about the extra information it provides beyond single-gene level p-values.

The tests based on PC modules treat highly correlated genes and their change between the groups as one unit. In this way, it points out that the transcriptional co-changes of certain genes are not accidental; they function as one group in the natural state and are expected to be impacted simultaneously under treatment or disease status. The anchored method also compares the predictive value across genes and selects the most discriminant ones. There may be many genes impacted by the treatment (and they are divided into correlation groups), and sparse classifiers typically select only a few from each group. It is more effective in pointing the researcher to the most major signal of the data set. The relevant interpretation discussion is revised in the ``Testing Results" section of the Lupus study.

% Kathryn said:

% We used the GLMNB for the Lupus analysis because we were originally writing for statisticians and it seemed more familiar for them.  Now writing for Biostatistics I’m afraid we have to use one of the major packages for DE analysis.  DESeq2 is famous. Tim Barry has his SCEPTRE package. I think the latter is better, but prob’y we should use the former.  It is also customary to use FDR instead of Bonf.

% \textcolor{teal}{For comparisons, and biological insights, we have a good coverage PC4.    We need to make a stronger statement about what PC1 is identifying from a biological viewpoint.   Maybe AI could help us here.  If you are careful you can get some good citations.}   I see you do have good citations for a lot of your other claims.

% They say we didn’t compare our results to the standard analysis (currently GLMNBadj), but actually we did.    Once we replace this with a DE package, then we need to see what we gained and lost.  I guess we simply try to write things with more punch to make whatever message we have stronger -- having citations always makes the claims more impactful.  Currently, our better message might be with PC1, where these genes are not significant with multiple testing correction.

% Again, I don’t know why they say we didn’t do biological interpretation b/c we did!  3 of the anchor genes are not predicted by other methods.  Is there anything we can say about why (statistically) they come up with this analysis and none of the other approaches? 

% \textcolor{teal}{Genes in a module are correlated.}

% \textcolor{teal}{PC1 contrasts GLM NB ADJ, it has more positive.}

\begin{shaded}
7. Justification of pseudo-bulk and normalization choices.
In Section 5.1, the authors aggregate counts to pseudo-bulk at the individual level and then apply shifted log-normalization. A brief rationale for this particular pipeline, along with comments on whether the proposed methods would still apply (or need modification) to cell-level data, would be useful—especially since different preprocessing choices are common in scRNA-seq practice.
\end{shaded}

The main motivation of aggregating single-cell level information to the individual level pseudo-bulk is to eliminate the dependency between cells within one individual. This choice didn't cause devastating power loss, as there is a good number of participants in this study. In the real-data section, we further clarified this motivation and acknowledged the difficulty of working with single-cell level information directly. The example in Section~5.2 is single-cell level, in which this individual-level clustering is absent. 

The proposed method in principle works for single-cell data without log normalization as it relies on the central limit theorem rather than a more restricted parametric assumption. We applied a log transformation to unify the scaling of different genes---and for many researchers this is considered a ``standard" pipeline. As a complement, the results in the second example were not log transformed (we removed the batch effect using Poisson regression).


\begin{shaded}
8. Multiple testing and interpretation across PCs.
For the debiased tests along PC1–PC4, it would be good to clarify how multiple testing is handled when jointly interpreting these projected nulls. Is any correction applied when simultaneously testing four projections? If not, perhaps note that the p-values are marginal for each direction and suggest how one might adjust them if using multiple PCs in practice.
\end{shaded}

Thank you for your careful reading. The reported p-values are marginal with no adjustment. They need to be adjusted to control the family-wise error rate when we conduct a large number of projective tests. We added this note to the results section of the Lupus study.

\newpage
\section{Reviewer 2's Comments}

\begin{shaded}
    This manuscript proposes adaptive projection-based two-sample tests for high-dimensional gene expression data.  The manuscript is generally clearly written and the methodological framework is conceptually appealing. However, several aspects of the paper require clarification and strengthening before the contribution can be fully assessed. Overall, I believe the paper contains promising ideas but requires substantial revision before it can be considered for publication.
\end{shaded}

\begin{shaded}
    1. Positioning relative to existing high-dimensional testing literature
    
The manuscript motivates projection-based testing using sparse principal component directions. However, there exists a substantial literature on projection-based and sparse high-dimensional two-sample testing procedures. The paper would benefit from a clearer discussion of how the proposed approach differs from or improves upon existing methods such as projection pursuit tests, sparse mean testing approaches, and screening-based two-sample testing procedures. Clarifying the conceptual novelty and positioning within the broader literature would significantly strengthen the contribution.
\end{shaded}

Thank you for the suggestions. We added more details in the introduction to highlight the contribution of this work:

\begin{itemize}
    \item For the one-step correction proposal, we work under the general theoretical framework. However, an eigenvector is not a simple linear mapping of the distribution. The exact details of the one-step correction, as well as its theoretical guarantee, were novel in this work. 
    \item For the anchored projection proposal, we emphasized that regular Gaussian statistics can be achieved under mild conditions without debiasing. Cross-fitting framework, using PC1 as the anchoring proxy and the equivalent means under the null hypothesis enabled this interesting regime. It has never been explored in the literature that this simple and application-relevant scheme would lead to regular inferential procedures. The closest one in the literature would be \cite{liu2022multiple}: However, their result was not established under the null hypothesis (because there was no proxy anchoring direction), the classifier aims at the unique optimal direction, and they implemented sample-splitting rather than cross-fitting (although they proposed procedures to enhance the power due to a smaller sample size).
\end{itemize}

We also added one paragraph at the end of the introduction section to discuss several canonical high-dimensional mean comparison methods. Although this work targets high-dimensional mean comparison, the specific mathematical formulation of our problem is novel and more structured compared with the most commonly considered global null. Our method was motivated more by the debiasing and cross-validation literature---so in this paragraph, we also emphasize this difference and try not to mislead the readers to the line of research reviewed in \cite{huang2022overview}.

\begin{shaded}
    2. Statistical formulation of the projected hypothesis

The projected null hypothesis is written as:
$$
H_{0, proj}(u): (\mu_X - \mu_Z)^\top u = 0
$$
where the projection direction $u$ is estimated from the data. When the projection direction is random, the null hypothesis effectively depends on a data-dependent quantity. While the manuscript proposes a debiasing approach to address this issue, it would be helpful to more clearly explain how the inferential validity is preserved when the projection direction is estimated from the same dataset used for testing. In particular, a clearer explanation of the theoretical justification for the proposed correction would improve readability.
\end{shaded}

Thank you for the suggestion. We revised the writing in the Intuition of One-step Correction section extensively. The current version better stresses the source of the irregular distribution of the plug-in estimator. It explains how the influence function appears naturally when people use a linear function to approximate the problematic term, which further motivates the bias-correction scheme.

\begin{shaded}
    3. Debiasing via double machine learning

The debiasing procedure based on the double machine learning framework is an interesting component of the proposed method. However, the manuscript provides only a relatively brief description of this step.

It would be helpful to clarify: 1. what the nuisance parameter is in this framework. 2. how sample splitting or cross-fitting is implemented in practice. 3. what assumptions are required to guarantee asymptotic validity of the test statistic.

Since DML-based inference relies on Neyman orthogonality and certain rate conditions for nuisance estimation, it would be useful to explicitly state the assumptions under which the proposed procedure is theoretically valid.
\end{shaded}

There are several quantities needed to be estimated from the left-out sample to perform the debiasing step. We managed to show, which is also one of the theoretical contributions, that all of them can be reduced to two items: the gene-gene covariance matrix $\Sigma$ and the mean expression level vector $\mu_X, \mu_Z$. This information is stated in the main theorem in section ``Debiased Tests for the Projected Null".

To clarify how the cross-fitting works, we included a self-contained description of the sample splitting strategy in the appendix section ``More Details on the Debiased Test", as you suggested in a later comment.

The assumptions needed to guarantee normality are stated/cited in the main theorem as well. We also expanded the discussion afterwards to include more methods and settings when these assumptions (essentially the $o_P(n^{-1/4})$ convergence rate) are satisfied. Now it discusses near-sparse true mean vectors as well as distributions with heavier tails.

\begin{shaded}
    4. Robustness of projection-based testing

Projection-based tests are generally most powerful when signals align with a small number of directions. It would be helpful if the authors could discuss how the proposed method performs when signals are distributed across many weak directions rather than concentrated in a sparse subspace. In addition, since the projection direction is estimated using PCA, the method implicitly assumes that the signal structure is reflected in the covariance structure. It would be useful to discuss how sensitive the method is to situations where mean differences occur in directions that are not strongly represented in the covariance eigenstructure.
\end{shaded}

When the signal is very much spread out or perfectly orthogonal to the space spanned by the leading eigenvectors, the proposed method should not be as powerful. Similar to many other methods reviewed in \cite{huang2022overview},  each high-dimensional comparison method has its own favorable scenarios and those it has more trouble with. This is also an expected feature of this proposal: when there is no signal aligned with $v_j$, it should not reject $(\mu_X - \mu_Z)^\top v_j = 0$. 

Beyond power considerations, controlling type-I error is an equally important goal when the signal is orthogonal to the PC directions. The plug-in test is appealing but cannot achieve this because it omits the variability in the estimated PC, and debiased testing can be a systematic solution. If one is not careful with the tools applied, the rejection rate may be much inflated even when the signal is perpendicular to the PC direction. 


\begin{shaded}
    5. Simulation design

The simulation study appears somewhat aligned with the structure of the proposed method. In particular, the covariance matrix used in the simulations is constructed so that the leading principal component corresponds to the signal direction.

Because the proposed method relies on identifying such directions, this design may favor projection-based tests. It would strengthen the empirical evaluation to include additional scenarios where: 1. the signal is not aligned with the leading principal component. 2. the signal is distributed across multiple directions. 3.the covariance structure is less structured.
\end{shaded}

In the appendix section ``Simulation Under Unstructured Sparse Correlation", we added some new simulation settings where the covariance matrix is less structured (but satisfies row-wise sparsity so that it is estimable). We observe that the debiased test controls the type-I error well and establishes increasing power when more samples are available.

We also highlight that when the signal is not aligned with the leading PC, as shown in the ``Projective Null" column, the proposed method is not expected to have a rejection rate $>5\%$. This is the null hypothesis. It is interestingly non-trivial to construct a method that can control type-I error under the projective null (in which the plug-in proposal failed).

\begin{shaded}
    6. Real data analysis

The real data example illustrates the application of the proposed method but could be strengthened by comparing the results with standard differential expression approaches commonly used in single-cell RNA-seq analysis. Providing some biological interpretation of the identified gene sets would also help demonstrate the practical utility of the method.
\end{shaded}

\begin{shaded}
Minor Comments

Page 3, Literature review.
The discussion of related work could be expanded to better position the proposed method relative to existing high-dimensional testing procedures.
\end{shaded}

We appreciate your suggestion. We added one paragraph at the end of the introduction section to discuss several canonical high-dimensional mean comparison methods. Although this work targets high-dimensional mean comparison, the specific mathematical formulation of our problem is novel and more structured compared with the most commonly considered global null. Our method was motivated more by the debiasing and cross-validation literature---so in this paragraph, we also emphasize this difference and try not to mislead the readers to the line of research reviewed in \cite{huang2022overview}.

\begin{shaded}
Page 5, Sparse PCA step.
The manuscript states that the projection direction is obtained via sparse PCA. However, the exact optimization problem used for sparse PCA is not specified. Providing the explicit objective function and describing how tuning parameters are selected would improve clarity and reproducibility.
\end{shaded}

We implemented the $\ell_1$-penalized sparse PCA framework introduced by \cite{witten2009penalized}. We added the following details in the appendix of the paper. 

\begin{quote}
The authors of \cite{witten2009penalized} applied an iterative algorithm to solve a penalized matrix decomposition of the data matrix, which can be viewed as a low-rank singular value decomposition with an $\ell_1$ constraint on the singular vectors. The explicit objective function is given in Equation~(2.7) of \cite{witten2009penalized}. The tuning parameter in this optimization problem is selected by cross-validation. The details of this procedure are described in the manual for the \texttt{SPC.cv} function in the \texttt{R} package \texttt{PMA} \cite{PMA}. We include a short summary here for readers' convenience.

The \texttt{SPC.cv} function performs cross-validation for the first sparse principal component. It repeats the following steps across folds: (1) replace a fraction of the entries in the data matrix with missing values; (2) apply \texttt{SPC} to this partially observed data matrix over a range of tuning parameter values, each time obtaining a rank-one approximation $udv^\top$ with sparse $v$; and (3) measure the mean squared prediction error for the missing entries created in step 1. The selected tuning parameter is the one with the lowest average mean squared prediction error in step 3.

To perform cross-validation for the second sparse principal component, the same procedure is applied to the deflated matrix $X-udv^\top$, where $udv^\top$ is the rank-one component obtained by applying \texttt{SPC} to the original data matrix $X$.
\end{quote}

\begin{shaded}
Page 6, Algorithm description.
It would be helpful to include a concise algorithmic summary describing the full testing pipeline, including projection estimation, sample splitting, debiasing, and p-value calculation.
\end{shaded}

Thank you for the suggestion. The manuscript provides a self-contained description for the pipeline in the renamed appendix section ``More Details on the Debiased Test".

\begin{shaded}
Page 8, Simulation dimensionality.
The simulations consider p = 100 variables. Many genomic datasets involve thousands of genes. Evaluating performance in larger dimensional settings would strengthen the empirical study.
\end{shaded}

We appreciate the suggestions. The intention of Figure 2 on this page is to show that even under mild dimensionality, the projective null does not induce normal statistics. The more extended simulations (including the original main simulation and the new supplementary ones) are conducted under a larger $p = 1000$. Many single-cell expression profiles measure many more genes, and we recommend trimming the list to those with the highest expression level or expression variability, 

\begin{shaded}
Supplementary Materials, Section A.
The covariance matrix used in the simulation is defined as:
$$
\Sigma = 3v_1 v_1^\top + I_p
$$
which ensures that the leading principal component corresponds to the signal direction. This design may favor projection-based tests. Additional simulation scenarios with more general covariance structures would help evaluate robustness.  
\end{shaded}

To verify the robustness of the proposal, a set of supplementary simulations was added to the appendix section ``Simulation Under Unstructured Sparse Correlation". In this study, we considered unstructured covariance matrices, which can be more general than the spike models. We also considered a different pipeline for estimating the covariance matrix and computing its PCs. The proposed pipeline still functions as expected even when many details are altered.

\bibliography{main}
\bibliographystyle{plain}
\end{document}
